\section{Hasil dan Pembahasan}

% Slide 19: Bab 4.1.1 - Akuisisi dan Transformasi Data Numerik
\begin{frame}{Bab 4.1.1 - Akuisisi dan Transformasi Data Numerik}
    \begin{itemize}
        \item \textbf{Data Sampel (BBCA):} Kenaikan harga dari 4515,50 ke 4869,32 (Jan 2021).
        \item \textbf{Kalkulasi Imbal Hasil:}
        \begin{itemize}
            \item \textbf{Simple Return ($R$):} $\frac{4869,32 - 4515,50}{4515,50} \approx 0,0783$
            \item \textbf{Log Return ($r$):} $\ln(1 + 0,0783) \approx 0,0754$
        \end{itemize}
        \item \textbf{Justifikasi Pemilihan:}
        \begin{itemize}
            \item $R$ digunakan untuk $\mu$ guna menjaga linearitas portofolio ($r_{true} \approx \hat{r}_p + \frac{1}{2}\text{Var}(R)$).
            \item $r$ digunakan untuk $\Sigma$ guna menjamin stasioneritas runtun waktu.
        \end{itemize}
    \end{itemize}
\end{frame}

% Slide 19b: Bab 4.1.2 - Matriks Kovariansi dan Risk Aversion
\begin{frame}{Bab 4.1.2 - Matriks Kovariansi dan Risk Aversion}
    \begin{itemize}
        \item \textbf{Matriks Kovariansi ($N=2$):}
        \begin{equation*}
            \Sigma_{N2} = \begin{pmatrix} 0,0171 & 0,0022 \\ 0,0022 & 0,0225 \end{pmatrix}
        \end{equation*}
        \item \textbf{Parameter Risk Aversion Endogen ($\gamma$):}
        \begin{itemize}
            \item Rasio Sharpe Agregat ($Z$): $\bar{\mu}_l / \bar{\sigma}_l = 0,1330 / 0,2221 \approx 0,5988$.
            \item Fungsi Sigmoid: $\gamma = 1 / (1 + e^Z) \approx \mathbf{0,3547}$.
        \end{itemize}
        \item \textbf{Interpretasi:} Nilai $\gamma$ rendah mencerminkan optimisme pasar pada jendela awal, memicu alokasi kapital yang lebih berani.
    \end{itemize}
\end{frame}

% Slide 20: Bab 4.2 - Formulasi EPG dan Hasil Nash
\begin{frame}{Bab 4.2 - Formulasi EPG dan Hasil Nash}
    \begin{itemize}
        \item \textbf{Insentif Marginal ($\Delta \Phi$):}
        \begin{equation*}
            \Delta \Phi(\vec{x}) = \frac{\mu_i}{k} - \frac{\gamma \sigma_i^2}{2k^2} - \gamma \sum_{j < i} \sigma_{ij} \frac{x_j}{k^2}
        \end{equation*}
        \item \textbf{Hasil Pencarian Nash Equilibrium ($N=2$):}
    \end{itemize}
    \begin{table}
        \footnotesize
        \centering
        \begin{tabular}{cc|c|c|}
        \cline{3-4}
        & & \multicolumn{2}{c|}{Aset TLKM ($x_2$)} \\ \cline{3-4}
        & & Status 0 & Status 1 \\ \hline
        \multicolumn{1}{|c|}{\multirow{2}{*}{Aset BBCA ($x_1$)}} & Status 0 & (0, 0) & (0, 0,1183) \\ \cline{2-4} 
        \multicolumn{1}{|c|}{} & Status 1 & \textbf{(0,1909, 0)} & (0,1901, 0,1175) \\ \hline
        \end{tabular}
        \caption{Nash Equilibrium tercapai pada $|10\rangle$ (Hanya BBCA terpilih).}
    \end{table}
\end{frame}

% Slide 21: Bab 4.3 - Derivasi EPG ke Hamiltonian Ising
\begin{frame}{Bab 4.3 - Derivasi EPG ke Hamiltonian Ising}
    \begin{itemize}
        \item \textbf{Langkah Transformasi:}
        \begin{enumerate}
            \item Definisikan fungsi energi: $E(\vec{x}) = -\Phi(\vec{x})$.
            \item Substitusi variabel biner ke spin: $x_i = (1 - s_i)/2$.
            \item Ekstraksi koefisien Pauli-Z ($\hat{Z}$).
        \end{enumerate}
        \item \textbf{Hasil Hamiltonian Numerik (Jan 2021):}
        \begin{equation*}
            \hat{\mathcal{H}}_{N2} \approx 0,0915 \hat{Z}_1 + 0,0533 \hat{Z}_2 + 0,5022 \hat{Z}_1 \hat{Z}_2 - 0,1470
        \end{equation*}
        \item \textbf{Analisis:} Koefisien interaksi didominasi oleh pinalti anggaran $A \approx 0,5$, menjamin validitas solusi kardinalitas.
    \end{itemize}
\end{frame}

% Slide 22: Bab 4.4 - Sirkuit Kuantum dan Gradien Energi
\begin{frame}{Bab 4.4 - Sirkuit Kuantum dan Gradien Energi}
    \begin{columns}
        \begin{column}{0.55\textwidth}
            \begin{itemize}
                \item \textbf{Konstruksi State:} $\ket{\psi(\boldsymbol{\theta})} = \hat{U}_{ent} \hat{U}_{rot}(\boldsymbol{\theta}) \ket{0\dots0}$.
                \item \textbf{Gradien Energi:} Dicari untuk meminimalkan $\langle \hat{\mathcal{H}} \rangle$:
                \begin{equation*}
                    \nabla_{\theta} E = \frac{\partial}{\partial \theta} \langle \psi(\theta) | \hat{H} | \psi(\theta) \rangle
                \end{equation*}
                \item \textbf{Sinergi:} Kombinasi gerbang rotasi $R_y, R_z$ dan keterbelitan CNOT membentuk lanskap energi yang dapat dilatih.
            \end{itemize}
        \end{column}
        \begin{column}{0.45\textwidth}
            \begin{figure}
                \centering
                \includegraphics[width=\textwidth]{Hasil_N2_GT/Analisis_Window_N2/2021-05-05_circuit.png}
                \caption{Arsitektur HEA sistem N=2.}
            \end{figure}
        \end{column}
    \end{columns}
\end{frame}

% Slide 23: Bab 4.5 - Penurunan Parameter Shift Rule (PSR)
\begin{frame}{Bab 4.5 - Penurunan Parameter Shift Rule (PSR)}
    \begin{itemize}
        \item \textbf{Aksioma:} Generator gerbang kuantum seringkali memiliki dua nilai eigen unik ($\pm r$).
        \item \textbf{Derivasi:} Ekspansi uniter gerbang rotasi $U(\theta) = e^{-i\theta G}$.
    \end{itemize}
    \begin{equation*}
        \frac{\partial E}{\partial \theta} = \frac{1}{2} \left[ E\left(\theta + \frac{\pi}{2}\right) - E\left(\theta - \frac{\pi}{2}\right) \right]
    \end{equation*}
    \begin{itemize}
        \item \textbf{Signifikansi:} Memberikan nilai gradien yang identik dengan turunan analitik tanpa limit $\epsilon \to 0$.
    \end{itemize}
\end{frame}

% Slide 24: Bab 4.6 - Derivasi dan Iterasi Numerik SPSA
\begin{frame}{Bab 4.6 - Derivasi dan Iterasi Numerik SPSA}
    \begin{itemize}
        \item \textbf{Ekspansi Taylor Orde-2:} SPSA membatalkan suku orde genap melalui selisih evaluasi maju-mundur ($E_+ - E_-$).
        \item \textbf{Estimator Gradien:} $\hat{g}_i = (E_+ - E_-) / (2 c_k \Delta_i)$.
        \item \textbf{Contoh Satu Langkah Iterasi:}
        \begin{itemize}
            \item \textbf{Input:} $\theta^{(0)} = [\pi/2, \pi/2]$, $\Delta_k = [1, -1]$.
            \item \textbf{Output:} $\Delta E = -0,0076 \rightarrow \hat{g} = [-0,038, 0,038]^T$.
            \item \textbf{Update:} $\theta^{(1)} = \theta^{(0)} - 0,1\hat{g} = [1,5746, 1,5670]^T$.
        \end{itemize}
    \end{itemize}
\end{frame}

% Slide 25: Bab 4.7 - Evolusi Entropi: Shannon ke Von Neumann
\begin{frame}{Bab 4.7 - Evolusi Entropi: Shannon ke Von Neumann}
    \begin{itemize}
        \item \textbf{Shannon ($H$):} Ketidakpastian informasi klasik $-\sum p_i \log_2 p_i$.
        \item \textbf{Von Neumann ($S$):} Perluasan ke operator densitas kuantum $\rho$:
    \end{itemize}
    \begin{equation*}
        S(\rho) = -\text{Tr}(\rho \log \rho) = -\sum_i \lambda_i \log \lambda_i
    \end{equation*}
    \begin{itemize}
        \item \textbf{Studi Kasus Jendela 1:}
        \begin{itemize}
            \item \textbf{N=2:} Probabilitas status optimal 0,98 $\rightarrow S = 0,1414$ bit.
            \item \textbf{N=4:} Probabilitas status optimal 0,94 $\rightarrow S = 0,3274$ bit.
        \end{itemize}
    \end{itemize}
\end{frame}

% Slide 26: Bab 4.8 - Analisis Barren Plateau dan Warm-Start
\begin{frame}{Bab 4.8 - Analisis Barren Plateau dan Warm-Start}
    \begin{itemize}
        \item \textbf{Jebakan Gradien:} Variansi gradien menyusut drastis pada sistem $N=4$ saat kedalaman $L > 5$.
        \item \textbf{Expressibility vs Trainability:} Peningkatan kapasitas sirkuit seringkali mengorbankan kemudahan optimasi.
    \end{itemize}
    \begin{center}
        \includegraphics[width=0.7\textwidth]{Hasil_N4_GT/grafik_depth_per_window_N4.png} \\
        \footnotesize \textit{Visualisasi pemilihan depth optimal menggunakan inisialisasi Nash.}
    \end{center}
    \begin{itemize}
        \item \textbf{Observasi:} Inisialisasi EPG memandu sirkuit ke sub-ruang Hilbert yang bermakna, menjaga gradien tetap terdeteksi.
    \end{itemize}
\end{frame}

% Slide 27: Bab 4.9 - Validasi Brute Force dan Benchmark SLSQP
\begin{frame}{Bab 4.9 - Validasi Brute Force dan Benchmark SLSQP}
    \begin{itemize}
        \item \textbf{Brute Force (N=2):}
        \begin{itemize}
            \item $E(|10\rangle) = -0,1909$ (Global Minimum).
            \item $E(|00\rangle) = 4,1972$ (Forbidden State).
        \end{itemize}
        \item \textbf{SLSQP:} Algoritma optimasi kontinu dengan diskritisasi \textit{Top-K}.
    \end{itemize}
    \begin{exampleblock}{Fair Comparison}
        Kedua metode diuji pada parameter batasan yang identik guna membandingkan efisiensi alokasi modal hibrida secara objektif.
    \end{exampleblock}
\end{frame}

% Slide 28: Bab 4.10 - Analisis Performa Portofolio (Hasil Utama)
\begin{frame}{Bab 4.10 - Analisis Performa Portofolio (Hasil Utama)}
    \begin{table}
        \scriptsize
        \centering
        \begin{tabular}{|l|c|c|c|c|c|c|}
        \hline
        \multirow{2}{*}{\textbf{Strategi}} & \multicolumn{3}{c|}{\textbf{Sistem N=2}} & \multicolumn{3}{c|}{\textbf{Sistem N=4}} \\ \cline{2-7} 
         & \textbf{Return} & \textbf{Sharpe} & \textbf{MDD} & \textbf{Return} & \textbf{Sharpe} & \textbf{MDD} \\ \hline
        \textbf{GT-VQE (Hibrida)} & \textbf{51,34\%} & \textbf{0,7357} & 15,91\% & \textbf{73,15\%} & \textbf{1,0107} & \textbf{16,68\%} \\ \hline
        VQE (Tanpa GT) & 43,21\% & 0,6471 & 15,91\% & 16,89\% & 0,3592 & 20,89\% \\ \hline
        SLSQP (Klasik) & 51,34\% & 0,6736 & 15,91\% & 12,91\% & 0,2810 & 21,12\% \\ \hline
        Equal Weight & 37,29\% & 0,6848 & 14,62\% & 35,57\% & 0,6790 & 18,42\% \\ \hline
        \end{tabular}
        \caption{GT-VQE mendominasi seluruh metrik pada dimensi sistem yang lebih besar (N=4).}
    \end{table}
\end{frame}

% Slide 30: Bab 4.11 - Dinamika Rebalancing Jendela Ekstrem
\begin{frame}{Bab 4.11 - Dinamika Rebalancing Jendela Ekstrem}
    \begin{itemize}
        \item \textbf{Agilitas Portofolio:} Transisi status $|10\rangle \to |01\rangle$ secara instan merespons lonjakan risiko BBCA (Feb-Mar 2021).
    \end{itemize}
    \begin{figure}
        \centering
        \begin{subfigure}[b]{0.45\textwidth}
            \centering
            \includegraphics[width=\textwidth]{Hasil_N2_GT/Analisis_Window_N2/2021-02-02_window.png}
            \caption{Jan-Feb 2021 (Aset BBCA)}
        \end{subfigure}
        \hfill
        \begin{subfigure}[b]{0.45\textwidth}
            \centering
            \includegraphics[width=\textwidth]{Hasil_N2_GT/Analisis_Window_N2/2021-03-04_window.png}
            \caption{Mar-Apr 2021 (Aset TLKM)}
        \end{subfigure}
        \caption{Quantum rebalancing mendeteksi pergeseran energi Hamiltonian pasar.}
    \end{figure}
\end{frame}
